\begin{frame}{Second question: can we do better?}

\small
\begin{columns}[T,totalwidth=\textwidth]
\begin{column}{0.36\textwidth}
\begin{block}{Idea}
The signal shape contains information about how the event was measured.
\end{block}

\begin{block}{Main question}
Can a model learn a waveform- dependent timing correction, rather than a black-box TOF prediction?
\end{block}

\begin{alertblock}{Risk with Machine Learning}
A very flexible model can make the time distribution narrower without necessarily learning the physical timing correction we want.
\end{alertblock}
\end{column}

\begin{column}{0.58\textwidth}
\centering

\smartimage[0.78\linewidth]{0.65\textheight}
{figures/distr_targ.png}
{Example from previous waveform-ML work.}

\vspace{0.15em}

{
Source: X.~Feng, A.~Muhashi, Y.~Onishi, R.~Ota, and H.~Liu,
``Transformer-CNN hybrid network for improving PET time of flight prediction,''
\textit{Physics in Medicine \& Biology}, 2024.
}
\end{column}
\end{columns}
\end{frame}
